% TeX root=../main.tex

\begin{table}[!htb]
	\caption{Sistema vocálico do indo-europeu ao \glsxtrlong{eslavoantigo}}\label{tab:sistvoc}
	\begin{center}
		\begin{tabular}[c]{lllll}
			\toprule
			\emph{Vogais}                                               & *ĭ > \emph{ь}                                          & *ĕ > \emph{e}                            & *ă, ŏ > *ă > \emph{o} & *ŭ > \emph{ъ} \\
			\emph{breves}                                               &                                                        & \hspace{1em} > \emph{o}\slash{}\emph{-u̯} &                       &               \\
			\midrule
			\emph{Vogais}                                               & *ī > \emph{i}                                          & *ē > \ocsex{E}
			[\ocsex{E}\textsubscript{1}]                                & *ā, *ō > *ā > \emph{a}                                 & *ū > \emph{y}                                                                    \\
			\emph{longas}                                               &                                                        &                                          &                       &               \\
			\midrule
			\emph{Ditongos}                                             &                                                        & *ei̯ > [\emph{i}\textsubscript{1}]        & *ai̯, *oi̯ > \ocsex{E}
			[\ocsex{E}\textsubscript{2}]/ \emph{i} [i\textsubscript{2}] &                                                                                                                                           \\
			                                                            &                                                        & *eu̯ > \ocsex{ju}                         & *au̯, *ou̯ > \emph{u}   &               \\
			                                                            & \multicolumn{2}{l}{*ī̆n, *ī̆m, }                         & *an, *am, *on, *om > \ocsex{õ}           &                                       \\
			                                                            & \multicolumn{2}{l}{\hspace{0.5em}*en, *em > \ocsex{ẽ}} & *ŭn, *ŭm > \ocsex{õ}                     &                                       \\
			\midrule
			\emph{Ressoantes}                                           & *r̥, *l̥ > \ocsex{rь lь}\slash\ocsex{rъ lъ}              &                                          &                       &               \\
			\emph{vocálicas}                                            & *n̥, *m̥ > \ocsex{ьnV}, \ocsex{ьmV}\slash{}\ocsex{ęC}    &                                          &                       &               \\
			\bottomrule
			\multicolumn{5}{l}{\small NB\@: após consoante palatalizada, *\pietrans{E} >
			\ocsex{E} > \emph{a}.}                                                                                                                                                                                  \\
		\end{tabular}
	\end{center}
\end{table}
